% ===========================================================================
%  Chapitre 10 — Calcul matriciel
%  PE 2026, composante ALGÈBRE
% ===========================================================================
\bmtchapitre{Calcul matriciel}
  {Effectuer et démontrer les calculs sur les matrices carrées d'ordre au plus
   quatre : somme, produit, déterminant, transposée, inverse, trace et
   puissances ; résoudre un système linéaire par les matrices.}
  {Algèbre \textbullet\ environ 10 heures}

Une matrice n'est rien d'autre qu'un tableau de nombres. Ce qui en fait un
objet mathématique, c'est qu'on décide de l'additionner et de la multiplier —
et que cette multiplication, contrairement à toutes celles que vous connaissez,
n'est pas commutative. C'est la première fois de votre scolarité que $AB$ et
$BA$ peuvent différer, et il faut quelques exercices pour s'y habituer.

L'intérêt du calcul matriciel apparaît dès qu'on écrit un système linéaire
sous la forme $AX = B$. Le système entier tient alors en une égalité, et le
résoudre revient à inverser une matrice. Le baccalauréat exploite cette idée
chaque année, presque toujours en la couplant à un raisonnement par récurrence
sur les puissances successives d'une matrice.

% ---------------------------------------------------------------------------
\begin{revision}
Rien de nouveau n'est requis : ces questions relèvent de la seconde et de la
première.

\begin{enumerate}[label=\textbf{\arabic*.}, leftmargin=*, itemsep=0.35em]
  \item Résoudre le système $\begin{cases} x+y = 5 \\ 2x-y = 1 \end{cases}$.
  \item Développer $(a+b)^{2}$ et $(a+b)(a-b)$.
  \item Rappeler le principe d'une démonstration par récurrence.
  \item Calculer $\begin{vmatrix} 2 & 3 \\ 1 & 4 \end{vmatrix}$, déterminant
        d'ordre deux vu en première.
  \item Qu'appelle-t-on un système linéaire ? Combien de solutions peut-il
        avoir ?
\end{enumerate}
\end{revision}

\begin{automatismes}
Sans calculatrice, sans brouillon.

\begin{enumerate}[label=\textbf{\alph*)}, leftmargin=*, itemsep=0.25em]
  \item $\matdeux{1}{2}{3}{4}+\matdeux{0}{1}{1}{0}$
  \item $2 \times \matdeux{1}{-1}{0}{3}$
  \item $\matdeux{1}{0}{0}{1} \times \matdeux{5}{7}{2}{3}$
  \item Déterminant de $\matdeux{3}{1}{2}{4}$
  \item Déterminant de $\matdeux{2}{4}{1}{2}$
  \item Transposée de $\matdeux{1}{2}{3}{4}$
  \item Trace de $\matdeux{5}{9}{4}{-2}$
  \item Carré de $\matdeux{0}{1}{0}{0}$
\end{enumerate}
\end{automatismes}

\begin{corrige}[automatismes]
a) $\matdeux{1}{3}{4}{4}$ \quad b) $\matdeux{2}{-2}{0}{6}$ \quad
c) $\matdeux{5}{7}{2}{3}$ \quad d) $10$ \quad e) $0$ \quad
f) $\matdeux{1}{3}{2}{4}$ \quad g) $3$ \quad h) la matrice nulle.
\end{corrige}

% ===========================================================================
\section{Définition d'une matrice carrée}

\begin{definition}[matrice carrée]
Une \textbf{matrice carrée d'ordre} $n$ est un tableau de $n$ lignes et $n$
colonnes de nombres réels. Le nombre situé à la ligne $i$ et à la colonne $j$
est le \textbf{coefficient} $a_{ij}$. On note $A = (a_{ij})$.

La \textbf{matrice identité} d'ordre $n$, notée $I_{n}$, a des $1$ sur la
diagonale et des $0$ partout ailleurs.
\end{definition}

\begin{remarque}
Deux matrices sont égales lorsqu'elles ont le même ordre et que tous leurs
coefficients coïncident, un par un. Il n'y a pas d'autre façon d'être égales :
on ne compare pas des matrices d'ordres différents.
\end{remarque}

\begin{definition}[matrices particulières]
Une matrice est \textbf{diagonale} si tous ses coefficients hors de la
diagonale sont nuls ; \textbf{triangulaire supérieure} si tous ceux situés
sous la diagonale sont nuls ; \textbf{symétrique} si $a_{ij} = a_{ji}$ pour
tous $i$ et $j$.
\end{definition}

% ===========================================================================
\section{Somme et produit ; non-commutativité du produit}

\begin{definition}[somme, produit par un réel]
La somme $A+B$ s'obtient en additionnant les coefficients de même position.
Le produit $\lambda A$ s'obtient en multipliant tous les coefficients par
$\lambda$.
\end{definition}

\begin{definition}[produit de deux matrices]
Le produit $AB$ est la matrice dont le coefficient de la ligne $i$ et de la
colonne $j$ vaut
\[
  (AB)_{ij} = \sum_{k=1}^{n} a_{ik}\,b_{kj},
\]
c'est-à-dire le produit terme à terme de la ligne $i$ de $A$ par la colonne
$j$ de $B$.
\end{definition}

\begin{visualisation}[ligne contre colonne]
\begin{center}
\begin{tikzpicture}[x=1cm, y=1cm, font=\footnotesize]
  \node at (0,0) {$\begin{pmatrix} \color{bmtLaterite}a_{11} & \color{bmtLaterite}a_{12} \\ a_{21} & a_{22} \end{pmatrix}$};
  \node at (2.2,0) {$\times$};
  \node at (4.3,0) {$\begin{pmatrix} b_{11} & \color{bmtVetiver}b_{12} \\ b_{21} & \color{bmtVetiver}b_{22} \end{pmatrix}$};
  \node at (6.4,0) {$=$};
  \node at (8.9,0) {$\begin{pmatrix} \ast & \color{bmtIndigo}c_{12} \\ \ast & \ast \end{pmatrix}$};
  \draw[bmtLaterite, line width=0.9pt, rounded corners=2pt]
    (-0.72,0.28) rectangle (0.6,-0.05);
  \draw[bmtVetiver, line width=0.9pt, rounded corners=2pt]
    (4.42,0.42) rectangle (4.9,-0.42);
  \node[bmtIndigo, anchor=north, align=center] at (8.9,-0.6)
    {$c_{12} = a_{11}b_{12}+a_{12}b_{22}$};
\end{tikzpicture}
\end{center}

Le coefficient qui s'inscrit à la ligne $i$ et à la colonne $j$ du résultat ne
dépend que de la ligne $i$ du facteur de gauche et de la colonne $j$ du facteur
de droite. Retenez le geste : on parcourt la ligne horizontalement et la
colonne verticalement, en même temps, et l'on additionne les produits
rencontrés.
\end{visualisation}

\begin{avertissement}
Le produit matriciel n'est pas commutatif : en général $AB \neq BA$. Deux
autres surprises l'accompagnent : un produit de deux matrices non nulles peut
être nul, et l'on ne peut pas simplifier une égalité $AB = AC$ par $A$. Toutes
les identités remarquables perdent donc leur validité ; en particulier
$(A+B)^{2} = A^{2}+AB+BA+B^{2}$, et cette expression ne se ramène à
$A^{2}+2AB+B^{2}$ que si $A$ et $B$ commutent.
\end{avertissement}

\begin{exemple}[deux produits qui diffèrent]
Prenons $A = \matdeux{0}{1}{0}{0}$ et $B = \matdeux{0}{0}{1}{0}$. Alors
\[
  AB = \matdeux{1}{0}{0}{0}
  \qquad\text{et}\qquad
  BA = \matdeux{0}{0}{0}{1}.
\]
Les deux produits diffèrent. On observe au passage que $A^{2}$ est la matrice
nulle sans que $A$ le soit : dans l'univers matriciel, un carré nul
n'entraîne pas la nullité.
\end{exemple}

\begin{entrainement}[somme et produit]
\exo[\niveau{1}] Calculer $A+B$ et $3A$ pour
$A = \matdeux{2}{-1}{0}{4}$ et $B = \matdeux{1}{1}{-2}{0}$.

\exo[\niveau{1}] Calculer $AB$ et $BA$ pour les matrices précédentes.

\exo[\niveau{2}] Soit $N = \mattrois{0}{1}{0}{0}{0}{1}{0}{0}{0}$. Calculer
$N^{2}$ et $N^{3}$.

\exo[\niveau{2}] Vérifier que $A = \matdeux{1}{1}{0}{1}$ et
$B = \matdeux{1}{2}{0}{1}$ commutent.

\exo[\niveau{3}] Déterminer toutes les matrices d'ordre deux qui commutent
avec $\matdeux{1}{1}{0}{1}$.
\end{entrainement}

% ===========================================================================
\section{Transposée et trace}

\begin{definition}[transposée, trace]
La \textbf{transposée} de $A$, notée $\transpo{A}$, est la matrice obtenue en
échangeant lignes et colonnes : son coefficient de position $(i\,;j)$ est
$a_{ji}$. La \textbf{trace} de $A$ est la somme de ses coefficients
diagonaux.
\end{definition}

\begin{propriete}[règles de calcul]
Pour toutes matrices carrées de même ordre,
\[
  \transpo{(A+B)} = \transpo{A}+\transpo{B}, \qquad
  \transpo{(AB)} = \transpo{B}\,\transpo{A}, \qquad
  \transpo{\left(\transpo{A}\right)} = A .
\]
La trace vérifie de plus $\operatorname{tr}(AB) = \operatorname{tr}(BA)$, même
lorsque $AB \neq BA$.
\end{propriete}

\begin{avertissement}
Dans la transposée d'un produit, l'ordre des facteurs s'inverse. Écrire
$\transpo{(AB)} = \transpo{A}\,\transpo{B}$ est une faute, et elle coûte
systématiquement des points.
\end{avertissement}

% ===========================================================================
\section{Déterminant d'une matrice carrée}

\begin{definition}[déterminant d'ordre deux et trois]
Pour $A = \matdeux{a}{b}{c}{d}$, on pose $\dete(A) = ad-bc$.

Pour une matrice d'ordre trois, le développement suivant la première ligne
donne
\[
  \dete
  \mattrois{a}{b}{c}{d}{e}{f}{g}{h}{i}
  = a\begin{vmatrix} e & f \\ h & i \end{vmatrix}
  - b\begin{vmatrix} d & f \\ g & i \end{vmatrix}
  + c\begin{vmatrix} d & e \\ g & h \end{vmatrix}.
\]
\end{definition}

\begin{remarque}
Les signes alternent selon le damier
$\begin{smallmatrix} + & - & + \\ - & + & - \\ + & - & + \end{smallmatrix}$.
On peut développer suivant n'importe quelle ligne ou colonne : on choisit
naturellement celle qui contient le plus de zéros. Pour une matrice
triangulaire, le déterminant est simplement le produit des coefficients
diagonaux.
\end{remarque}

\begin{propriete}[propriétés du déterminant]
$\dete(AB) = \dete(A)\times\dete(B)$ et
$\dete\!\left(\transpo{A}\right) = \dete(A)$. En revanche,
$\dete(A+B)$ n'a aucune raison de valoir $\dete(A)+\dete(B)$.
\end{propriete}

\begin{entrainement}[déterminants]
\exo[\niveau{1}] Calculer le déterminant de $\matdeux{5}{2}{3}{4}$ et de
$\matdeux{6}{-2}{-3}{1}$.

\exo[\niveau{2}] Calculer le déterminant de
$\mattrois{1}{0}{2}{0}{1}{3}{0}{0}{1}$ et de
$\mattrois{2}{3}{5}{1}{0}{2}{0}{0}{-1}$.

\exo[\niveau{2}] Pour quelles valeurs de $m$ la matrice
$\matdeux{m}{1}{4}{m}$ a-t-elle un déterminant nul ?

\exo[\niveau{3}] Montrer que si $A$ est d'ordre $3$, alors
$\dete(2A) = 8\dete(A)$.
\end{entrainement}

% ===========================================================================
\section{Matrice inverse : existence et calcul}

\begin{definition}[matrice inversible]
Une matrice carrée $A$ d'ordre $n$ est \textbf{inversible} s'il existe une
matrice $B$ de même ordre telle que
\[
  AB = BA = I_{n}.
\]
Cette matrice $B$, lorsqu'elle existe, est unique ; on la note $A^{-1}$.
\end{definition}

\begin{theoreme}[critère d'inversibilité]
$A$ est inversible si et seulement si $\dete(A) \neq 0$. Pour une matrice
d'ordre deux,
\[
  \matdeux{a}{b}{c}{d}^{-1}
  = \frac{1}{ad-bc}\matdeux{d}{-b}{-c}{a}.
\]
\end{theoreme}

\begin{demonstration}
Si $A$ est inversible, $\dete(A)\dete\!\left(A^{-1}\right) = \dete(I_{n}) = 1$,
donc $\dete(A) \neq 0$. Réciproquement, pour l'ordre deux, on vérifie
directement le produit :
\[
  \matdeux{a}{b}{c}{d}\matdeux{d}{-b}{-c}{a}
  = \matdeux{ad-bc}{0}{0}{ad-bc} = (ad-bc)\,I_{2},
\]
et il suffit de diviser par $ad-bc$, non nul.
\end{demonstration}

\begin{methode}{trouver l'inverse à partir d'une relation polynomiale}
C'est la méthode attendue au baccalauréat, et elle évite tout calcul lourd.
Lorsque l'énoncé fournit une égalité du type
\[
  A^{2}+\alpha A+\beta I_{n} = 0 \quad\text{avec}\quad \beta \neq 0,
\]
on factorise par $A$ :
$A\left(A+\alpha I_{n}\right) = -\beta I_{n}$, puis on divise par $-\beta$ :
\[
  A \times \left(-\frac{1}{\beta}\right)\left(A+\alpha I_{n}\right) = I_{n},
\]
ce qui donne à la fois l'inversibilité et l'inverse. Aucune inversion directe
n'est nécessaire.
\end{methode}

\begin{exemple}[la méthode déroulée]
Soit $A$ vérifiant $2A-A^{2} = I_{3}$. En factorisant,
$A\left(2I_{3}-A\right) = I_{3}$ : la matrice $A$ est inversible et son
inverse est $P = 2I_{3}-A$. La vérification tient en une ligne, et l'on
n'a jamais eu besoin des coefficients de $A$.
\end{exemple}

\begin{entrainement}[inverse]
\exo[\niveau{1}] Déterminer l'inverse de $\matdeux{3}{1}{5}{2}$.

\exo[\niveau{2}] Soit $A = \mattrois{1}{0}{2}{0}{1}{3}{0}{0}{1}$. Vérifier que
$2A-A^{2} = I_{3}$ et en déduire $A^{-1}$.

\exo[\niveau{2}] Soit $A$ telle que $A^{3} = I_{3}$. Montrer que $A$ est
inversible et donner $A^{-1}$.

\exo[\niveau{3}] Soit $A$ vérifiant $4A+A^{2}-A^{3} = 4I_{3}$. En déduire
$A^{-1}$ en fonction de $A$ et $A^{2}$.
\end{entrainement}

% ===========================================================================
\section{Puissances successives d'une matrice}

\begin{definition}[puissance]
Pour $n \in \N^{*}$, $A^{n} = \underbrace{A \times A \times \cdots \times A}_{n
\text{ facteurs}}$, et l'on convient que $A^{0} = I_{n}$.
\end{definition}

\begin{methode}{calculer $A^{n}$}
Trois chemins, à choisir selon ce que fournit l'énoncé.
\begin{enumerate}[label=\textbf{\arabic*.}, leftmargin=*, itemsep=0.15em]
  \item \emph{Conjecture puis récurrence.} On calcule $A^{2}$ et $A^{3}$, on
        devine la forme générale, on la démontre par récurrence. C'est la voie
        la plus fréquente au baccalauréat.
  \item \emph{Décomposition $A = \lambda I+N$ avec $N$ nilpotente.} La formule
        du binôme s'applique car $I$ commute avec tout, et la somme s'arrête
        dès que $N$ s'annule.
  \item \emph{Diagonalisation fournie.} Si l'énoncé donne $P$ telle que
        $A = PBP^{-1}$ avec $B$ diagonale, alors $A^{n} = PB^{n}P^{-1}$, et
        $B^{n}$ se lit immédiatement.
\end{enumerate}
\end{methode}

\begin{exemple}[une récurrence type]
Soit $C = \mattrois{2}{2}{0}{0}{2}{0}{0}{0}{2}$. On calcule
$C^{2} = \mattrois{4}{8}{0}{0}{4}{0}{0}{0}{4}$, ce qui suggère
\[
  C^{n} = \mattrois{2^{n}}{n2^{n}}{0}{0}{2^{n}}{0}{0}{0}{2^{n}}.
\]
\emph{Initialisation.} Pour $n = 1$ la formule redonne $C$.

\emph{Hérédité.} Supposons-la vraie au rang $n$. Alors $C^{n+1} = C^{n}C$, et
le coefficient de position $(1\,;2)$ vaut
$2^{n}\times2+n2^{n}\times2 = (n+1)2^{n+1}$, tandis que les coefficients
diagonaux valent $2^{n}\times2 = 2^{n+1}$ : c'est bien la formule au rang
$n+1$.

\emph{Conclusion.} La formule est vraie pour tout entier $n \geq 1$.
\end{exemple}

\begin{entrainement}[puissances]
\exo[\niveau{1}] Calculer $A^{2}$ et $A^{3}$ pour
$A = \matdeux{1}{1}{0}{1}$, puis conjecturer $A^{n}$.

\exo[\niveau{2}] Démontrer la conjecture précédente par récurrence.

\exo[\niveau{2}] Soit $B = \mattrois{2}{0}{0}{0}{3}{0}{0}{0}{6}$. Donner
$B^{n}$.

\exo[\niveau{3}] Soit $A = PBP^{-1}$ avec $B$ diagonale. Démontrer par
récurrence que $A^{n} = PB^{n}P^{-1}$.
\end{entrainement}

% ===========================================================================
\section{Résolution d'un système linéaire}

\begin{propriete}[écriture matricielle]
Un système de $n$ équations linéaires à $n$ inconnues s'écrit $AX = B$, où
$A$ est la matrice des coefficients, $X$ la colonne des inconnues et $B$ la
colonne des seconds membres. Si $A$ est inversible, le système admet une
solution unique
\[
  X = A^{-1}B .
\]
\end{propriete}

\begin{demonstration}
En multipliant $AX = B$ à gauche par $A^{-1}$, on obtient $X = A^{-1}B$ ;
réciproquement cette colonne vérifie $A\left(A^{-1}B\right) = B$.
\end{demonstration}

\begin{avertissement}
La multiplication se fait \emph{à gauche} des deux côtés. Écrire
$X = BA^{-1}$ n'a pas de sens ici : les formats ne s'accordent même pas.
\end{avertissement}

\begin{exemple}[un système résolu par les matrices]
Résolvons
\[
  \begin{cases}
    3y+2z = 7 \\ x+y-z = 2 \\ 2x+y-3z = 1
  \end{cases}
\]
On pose $A = \mattrois{0}{3}{2}{1}{1}{-1}{2}{1}{-3}$,
$X = \begin{pmatrix} x \\ y \\ z \end{pmatrix}$ et
$B = \begin{pmatrix} 7 \\ 2 \\ 1 \end{pmatrix}$, de sorte que le système
s'écrit $AX = B$.

L'énoncé fournit $P = \mattrois{-2}{11}{-5}{1}{-4}{2}{-1}{6}{-3}$ ; le calcul
donne $AP = PA = I_{3}$, donc $P = A^{-1}$. Il vient
\[
  X = PB = \begin{pmatrix} -14+22-5 \\ 7-8+2 \\ -7+12-3 \end{pmatrix}
    = \begin{pmatrix} 3 \\ 1 \\ 2 \end{pmatrix}.
\]
On vérifie : $3+4 = 7$, $3+1-2 = 2$, $6+1-6 = 1$.
\end{exemple}

\begin{entrainement}[systèmes]
\exo[\niveau{1}] Écrire sous forme matricielle le système
$\begin{cases} 2x+y = 3 \\ x-y = 0 \end{cases}$, puis le résoudre.

\exo[\niveau{2}] Résoudre par les matrices
$\begin{cases} x+2y = 1 \\ 3x+5y = 4 \end{cases}$.

\exo[\niveau{3}] Discuter, suivant $m$, le nombre de solutions du système
$\begin{cases} mx+y = 1 \\ x+my = 1 \end{cases}$.
\end{entrainement}

% ===========================================================================
\begin{annales}
Le calcul matriciel occupe presque toujours la première partie de l'exercice 1
de l'épreuve, et le scénario varie peu : une relation polynomiale conduit à
l'inverse, puis une récurrence donne les puissances. Les énoncés qui suivent
sont repris de sujets réellement posés.

\exoann{Baccalauréat, Madagascar, série S, session 2026 — exercice 1, partie I}
On considère les matrices
\[
  A = \mattrois{1}{0}{2}{0}{1}{3}{0}{0}{1}, \qquad
  B = \mattrois{2}{3}{5}{1}{0}{2}{0}{0}{-1}, \qquad I_{3}.
\]
\begin{enumerate}[label=\textbf{\arabic*)}, leftmargin=*, itemsep=0.15em]
  \item Calculer le déterminant de $A$ et conclure.
  \item Vérifier que $2A-A^{2} = I_{3}$, puis en déduire la matrice $P$ telle
        que $A \times P = I_{3}$.
\end{enumerate}

\exoann{Baccalauréat, Madagascar, série S, session 2026 — exercice 1, partie I, suite}
On pose $D = ABP$, où $P$ est la matrice trouvée à l'exercice précédent.
Démontrer par récurrence que, pour tout $n \in \N^{*}$,
$D^{n} = AB^{n}P$.

\exoann{Baccalauréat, Madagascar, série S, session 2021 — exercice 1, partie I}
On considère
\[
  A = \mattrois{0}{1}{2}{1}{3}{0}{1}{2}{1}, \qquad
  B = \frac13\mattrois{-3}{-3}{6}{1}{2}{-2}{1}{-1}{1}, \qquad
  C = \mattrois{2}{2}{0}{0}{2}{0}{0}{0}{2}.
\]
\begin{enumerate}[label=\textbf{\arabic*)}, leftmargin=*, itemsep=0.15em]
  \item Montrer que $B$ est l'inverse de $A$.
  \item Calculer $A+C$.
\end{enumerate}

\exoann{Baccalauréat, Madagascar, série S, session 2021 — exercice 1, partie I, suite}
Avec la matrice $C$ de l'exercice précédent, démontrer par récurrence que,
pour tout $n \in \N^{*}\setminus\{1\}$,
\[
  C^{n} = \mattrois{2^{n}}{n2^{n}}{0}{0}{2^{n}}{0}{0}{0}{2^{n}}.
\]

\exoann{Baccalauréat, Madagascar, série S, session 2022 — exercice 1, partie II}
On considère la matrice
$A = \mattrois{0}{-1}{-1}{-1}{0}{-1}{-1}{-1}{0}$.
\begin{enumerate}[label=\textbf{\arabic*)}, leftmargin=*, itemsep=0.15em]
  \item Calculer $A^{2}$ et $A^{3}$.
  \item Vérifier que $4A+A^{2}-A^{3} = 4I_{3}$.
  \item En déduire que $A$ est inversible et donner $A^{-1}$.
\end{enumerate}

\exoann{Baccalauréat, Madagascar, série S, session 2023 — exercice 1, partie A}
On considère le système
$(S) : \begin{cases} 3y+2z = 7 \\ x+y-z = 2 \\ 2x+y-3z = 1 \end{cases}$
\begin{enumerate}[label=\textbf{\arabic*)}, leftmargin=*, itemsep=0.15em]
  \item Écrire $(S)$ sous la forme matricielle $AX = B$, où $A$ est carrée
        d'ordre trois.
  \item Soit $P = \mattrois{-2}{11}{-5}{1}{-4}{2}{-1}{6}{-3}$. Calculer
        $A \times P$ et $P \times A$, puis conclure.
  \item En déduire la matrice $X$.
\end{enumerate}

\exoann{Baccalauréat blanc, lycée Philibert Tsiranana, Terminale S, 2021-2022 — exercice 1, partie I}
On considère les matrices carrées d'ordre trois
\[
  A = \mattrois{3}{1}{1}{1}{5}{1}{1}{1}{3}, \qquad
  B = \mattrois{2}{0}{0}{0}{3}{0}{0}{0}{6},
\]
\[
  P = \mattrois{1}{-1}{1}{0}{1}{2}{-1}{-1}{1}, \qquad
  Q = \mattrois{3}{0}{-3}{-2}{2}{-2}{1}{2}{1}.
\]
\begin{enumerate}[label=\textbf{\alph*)}, leftmargin=*, itemsep=0.15em]
  \item Calculer $P \times Q$ et $B^{2}$.
  \item En déduire la matrice inverse $P^{-1}$ et la matrice $B^{n}$ pour
        $n \in \N^{*}$.
\end{enumerate}

\exoann{Baccalauréat blanc, lycée Philibert Tsiranana, Terminale S, 2021-2022 — exercice 1, partie I, suite}
Avec les matrices précédentes :
\begin{enumerate}[label=\textbf{\alph*)}, leftmargin=*, itemsep=0.15em]
  \item Montrer que $A = P\,B\,P^{-1}$.
  \item En déduire par récurrence que $A^{n} = P\,B^{n}\,P^{-1}$ pour tout
        $n \in \N^{*}$.
\end{enumerate}
\end{annales}

\begin{corrige}[annales]
\textbf{1.} La matrice $A$ est triangulaire supérieure : son déterminant est
le produit des coefficients diagonaux, soit $1$. Il est non nul, donc $A$ est
inversible. On calcule $A^{2} = \mattrois{1}{0}{4}{0}{1}{6}{0}{0}{1}$, d'où
$2A-A^{2} = I_{3}$. En factorisant, $A\left(2I_{3}-A\right) = I_{3}$ :
\[
  P = 2I_{3}-A = \mattrois{1}{0}{-2}{0}{1}{-3}{0}{0}{1}.
\]

\textbf{2.} \emph{Initialisation.} $D^{1} = ABP$ par définition.
\emph{Hérédité.} Si $D^{n} = AB^{n}P$, alors
$D^{n+1} = D^{n}D = AB^{n}P \times ABP = AB^{n}\left(PA\right)BP$. Or
$PA = I_{3}$ puisque $AP = I_{3}$ et que l'inverse est bilatère : il reste
$D^{n+1} = AB^{n+1}P$.

\textbf{3.} Le produit $AB$ se calcule ligne par colonne et donne $I_{3}$ ; il
en va de même de $BA$ : $B$ est bien l'inverse de $A$. Ensuite
$A+C = \mattrois{2}{3}{2}{1}{5}{0}{1}{2}{3}$.

\textbf{4.} \emph{Initialisation.} $C^{2} = \mattrois{4}{8}{0}{0}{4}{0}{0}{0}{4}$,
ce qui est la formule pour $n = 2$. \emph{Hérédité.} En multipliant la forme
supposée par $C$, le coefficient $(1\,;2)$ devient
$2^{n}\times2+n2^{n}\times2 = (n+1)2^{n+1}$ et les coefficients diagonaux
$2^{n+1}$ : c'est la formule au rang $n+1$.

\textbf{5.} En posant $J$ la matrice dont tous les coefficients valent $1$, on
a $A = I_{3}-J$ et $J^{2} = 3J$. Donc
$A^{2} = I_{3}+J$ et $A^{3} = I_{3}-3J$. Il vient
$4A+A^{2}-A^{3} = 4I_{3}-4J+I_{3}+J-I_{3}+3J = 4I_{3}$. En factorisant,
$A\left(4I_{3}+A-A^{2}\right) = 4I_{3}$, donc $A$ est inversible et
\[
  A^{-1} = \tfrac14\left(4I_{3}+A-A^{2}\right) = I_{3}-\tfrac12 J
  = \tfrac12\mattrois{1}{-1}{-1}{-1}{1}{-1}{-1}{-1}{1}.
\]

\textbf{6.} On pose $A = \mattrois{0}{3}{2}{1}{1}{-1}{2}{1}{-3}$,
$X = \transpo{(x\ y\ z)}$ et $B = \transpo{(7\ 2\ 1)}$. Le calcul donne
$AP = PA = I_{3}$, donc $P = A^{-1}$ et $X = PB$, soit $x = 3$, $y = 1$,
$z = 2$. La vérification dans les trois équations est immédiate.

\textbf{7.} Le produit $PQ$ vaut $6I_{3}$, donc $P^{-1} = \frac16 Q$. La
matrice $B$ étant diagonale, $B^{2} = \mattrois{4}{0}{0}{0}{9}{0}{0}{0}{36}$
et plus généralement
$B^{n} = \mattrois{2^{n}}{0}{0}{0}{3^{n}}{0}{0}{0}{6^{n}}$.

\textbf{8.} Le produit $PBP^{-1}$ se calcule et redonne $A$. Pour la
récurrence : la formule est vraie au rang $1$ ; si $A^{n} = PB^{n}P^{-1}$,
alors
$A^{n+1} = PB^{n}P^{-1}\times PBP^{-1} = PB^{n}\left(P^{-1}P\right)BP^{-1}
= PB^{n+1}P^{-1}$.
\end{corrige}

% ===========================================================================
\begin{jecherche}[les trois ateliers d'une menuiserie]
Une menuiserie fabrique des tabourets, des chaises et des tables. Chaque
meuble passe par trois ateliers : découpe, assemblage, finition. Le tableau
donne le nombre d'heures nécessaires dans chaque atelier, pour chaque meuble.

\begin{center}
\setlength{\tabcolsep}{9pt}
\renewcommand{\arraystretch}{1.2}
\begin{tabular}{@{}lccc@{}}
\toprule
& \textsf{\footnotesize\bfseries TABOURET} & \textsf{\footnotesize\bfseries CHAISE}
& \textsf{\footnotesize\bfseries TABLE} \\
\midrule
Découpe    & 1 & 2 & 3 \\
Assemblage & 1 & 3 & 4 \\
Finition   & 1 & 1 & 1 \\
\bottomrule
\end{tabular}
\end{center}

Une semaine donnée, les trois ateliers ont travaillé respectivement $30$,
$41$ et $16$ heures.

\begin{enumerate}[label=\textbf{\arabic*.}, leftmargin=*, itemsep=0.3em]
  \item Traduire la situation par un système de trois équations, puis par une
        égalité matricielle $AX = B$.
  \item Calculer $\dete(A)$ et conclure quant au nombre de productions
        compatibles avec un relevé donné.
  \item Déterminer $A^{-1}$.
  \item En déduire le nombre de meubles de chaque sorte fabriqués cette
        semaine.
  \item La semaine suivante, les relevés sont $24$, $33$ et $13$ heures.
        Donner la production correspondante.
  \item L'atelier de finition change de méthode : il lui faut désormais
        $2$~heures pour une chaise et $3$~heures pour une table, la ligne
        « finition » devenant $(1\,;2\,;3)$. Que devient le déterminant, et
        qu'en conclure sur la lecture des relevés ?
\end{enumerate}
\end{jecherche}

\begin{corrige}[les trois ateliers d'une menuiserie]
\textbf{1.} En notant $x$, $y$, $z$ les nombres de tabourets, de chaises et de
tables :
\[
  \begin{cases} x+2y+3z = 30 \\ x+3y+4z = 41 \\ x+y+z = 16 \end{cases}
  \qquad\text{soit}\qquad
  \mattrois{1}{2}{3}{1}{3}{4}{1}{1}{1}
  \begin{pmatrix} x \\ y \\ z \end{pmatrix}
  = \begin{pmatrix} 30 \\ 41 \\ 16 \end{pmatrix}.
\]

\textbf{2.} En développant suivant la première ligne,
$\dete(A) = 1(3-4)-2(1-4)+3(1-3) = -1+6-6 = -1$. Le déterminant étant non
nul, $A$ est inversible : à chaque relevé correspond une production et une
seule.

\textbf{3.} Le calcul des cofacteurs donne
\[
  A^{-1} = \mattrois{1}{-1}{1}{-3}{2}{1}{2}{-1}{-1},
\]
ce que confirme la vérification $AA^{-1} = I_{3}$.

\textbf{4.} $X = A^{-1}B$ donne
$x = 30-41+16 = 5$, $y = -90+82+16 = 8$ et $z = 60-41-16 = 3$ : cinq
tabourets, huit chaises et trois tables. Les trois équations sont bien
vérifiées.

\textbf{5.} Le même calcul avec $B = \transpo{(24\ 33\ 13)}$ donne
$x = 4$, $y = 7$, $z = 2$.

\textbf{6.} La ligne « finition » devient identique à la ligne « découpe ».
Deux lignes égales rendent le déterminant nul : la matrice n'est plus
inversible. Un relevé ne détermine alors plus la production — soit il est
impossible, soit il correspond à une infinité de productions. Autrement dit,
mesurer deux fois la même chose ne renseigne pas davantage.
\end{corrige}

% ===========================================================================
\begin{jemeteste}
Répondre par vrai ou faux, en justifiant en une ligne.

\begin{enumerate}[label=\textbf{\arabic*.}, leftmargin=*, itemsep=0.28em]
  \item Pour toutes matrices carrées, $AB = BA$.
  \item Si $A^{2} = 0$ alors $A = 0$.
  \item Si $\dete(A) = 0$, le système $AX = B$ n'a jamais de solution.
  \item $\transpo{(AB)} = \transpo{B}\,\transpo{A}$.
  \item Une matrice triangulaire est inversible si et seulement si aucun
        coefficient diagonal n'est nul.
  \item $\dete(A+B) = \dete(A)+\dete(B)$.
\end{enumerate}
\end{jemeteste}

\begin{corrige}[je me teste]
\textbf{1.} Faux — le produit matriciel n'est pas commutatif.
\textbf{2.} Faux — $\matdeux{0}{1}{0}{0}$ est un contre-exemple.
\textbf{3.} Faux — il peut y avoir une infinité de solutions ; ce qui est
perdu, c'est l'unicité.
\textbf{4.} Vrai.
\textbf{5.} Vrai — le déterminant est le produit des coefficients diagonaux.
\textbf{6.} Faux — prendre $A = I_{2}$ et $B = -I_{2}$ : le membre de gauche
vaut $0$, celui de droite $2$.
\end{corrige}

% ===========================================================================
\begin{commeaubac}
\bmtsource{Baccalauréat, série S, session 2026 — exercice 1, partie I}
On considère les matrices carrées
\[
  A = \mattrois{1}{0}{2}{0}{1}{3}{0}{0}{1}, \qquad
  B = \mattrois{2}{3}{5}{1}{0}{2}{0}{0}{-1}, \qquad I_{3}.
\]
\begin{enumerate}[label=\textbf{\arabic*.}, leftmargin=*, itemsep=0.25em]
  \item \begin{enumerate}[label=\textbf{\alph*)}, leftmargin=1.4em]
          \item Calculer le déterminant de $A$. Conclure.
                \hfill\textup{(0,5 pt)}
          \item Vérifier que $2A-A^{2} = I_{3}$. En déduire la matrice $P$
                telle que $A \times P = I_{3}$. \hfill\textup{(0,5 pt)}
        \end{enumerate}
  \item On pose $D = ABP$. Montrer par récurrence que, pour tout
        $n \in \N^{*}$, $D^{n} = AB^{n}P$. \hfill\textup{(0,5 pt)}
\end{enumerate}
\end{commeaubac}

\begin{corrige}[comme au baccalauréat]
\textbf{1. a)} $A$ est triangulaire, donc $\dete(A) = 1 \neq 0$ : elle est
inversible.
\textbf{b)} $A^{2} = \mattrois{1}{0}{4}{0}{1}{6}{0}{0}{1}$, donc
$2A-A^{2} = I_{3}$ et $A\left(2I_{3}-A\right) = I_{3}$ :
$P = \mattrois{1}{0}{-2}{0}{1}{-3}{0}{0}{1}$.

\textbf{2.} La formule est vraie au rang $1$. Si $D^{n} = AB^{n}P$, alors
$D^{n+1} = AB^{n}P\,ABP$ ; comme $PA = I_{3}$, il reste $AB^{n+1}P$. La
propriété est héréditaire, donc vraie pour tout $n \geq 1$.
\end{corrige}

% ===========================================================================
\begin{concours}
\bmtsource{D'après un concours d'entrée en école d'ingénieurs}
Soit $J$ la matrice carrée d'ordre $3$ dont tous les coefficients valent $1$,
et soit $A = aI_{3}+bJ$, où $a$ et $b$ sont deux réels.

\begin{enumerate}[label=\textbf{\arabic*.}, leftmargin=*, itemsep=0.25em]
  \item Calculer $J^{2}$, puis $A^{2}$ en fonction de $I_{3}$ et de $J$.
  \item Démontrer par récurrence qu'il existe deux suites $(\alpha_{n})$ et
        $(\beta_{n})$ telles que $A^{n} = \alpha_{n}I_{3}+\beta_{n}J$, et
        exprimer $\alpha_{n+1}$ et $\beta_{n+1}$ en fonction de
        $\alpha_{n}$ et $\beta_{n}$.
  \item En déduire une condition nécessaire et suffisante sur $a$ et $b$ pour
        que $A$ soit inversible, et donner alors $A^{-1}$ sous la forme
        $\lambda I_{3}+\mu J$.
\end{enumerate}
\end{concours}

\begin{corrige}[vers le concours]
\textbf{1.} Chaque coefficient de $J^{2}$ est la somme de trois produits de
$1$ : $J^{2} = 3J$. Alors
$A^{2} = a^{2}I_{3}+2abJ+b^{2}J^{2} = a^{2}I_{3}+\left(2ab+3b^{2}\right)J$.

\textbf{2.} La forme est vraie pour $n = 1$ avec $\alpha_{1} = a$ et
$\beta_{1} = b$. Si $A^{n} = \alpha_{n}I_{3}+\beta_{n}J$, alors
\[
  A^{n+1} = \left(\alpha_{n}I_{3}+\beta_{n}J\right)\left(aI_{3}+bJ\right)
  = a\alpha_{n}I_{3}
  +\left(b\alpha_{n}+a\beta_{n}+3b\beta_{n}\right)J .
\]
Donc $\alpha_{n+1} = a\alpha_{n}$ et
$\beta_{n+1} = b\alpha_{n}+(a+3b)\beta_{n}$.

\textbf{3.} Cherchons $A^{-1} = \lambda I_{3}+\mu J$. Le produit vaut
$a\lambda I_{3}+\left(b\lambda+a\mu+3b\mu\right)J$, et l'on veut
$I_{3}$ : il faut $a\lambda = 1$ et $b\lambda+(a+3b)\mu = 0$. Le système est
résoluble si et seulement si $a \neq 0$ et $a+3b \neq 0$, et l'on obtient
alors
\[
  \lambda = \frac1a, \qquad \mu = \frac{-b}{a(a+3b)},
  \qquad\text{soit}\qquad
  A^{-1} = \frac1a I_{3}-\frac{b}{a(a+3b)}J .
\]
\end{corrige}

% ===========================================================================
\begin{notepourlaclasse}
Consacrez la première séance entière à la non-commutativité. Tant que les
élèves n'ont pas eux-mêmes produit deux matrices dont les produits diffèrent,
ils continuent d'appliquer les identités remarquables, et toutes les
démonstrations du chapitre s'effondrent.

La méthode de l'inverse par relation polynomiale doit être vue comme la
méthode normale, non comme une astuce : c'est celle que le sujet impose chaque
année. Faites-la reconnaître sous ses différents habillages —
$2A-A^{2} = I$, $4A+A^{2}-A^{3} = 4I$, $A^{3} = I$ — jusqu'à ce que le geste
de factorisation par $A$ devienne réflexe.

Les récurrences matricielles échouent presque toujours au même endroit :
l'élève écrit $A^{n+1} = A^{n}A$ puis renonce à effectuer le produit. Exigez
le calcul explicite d'au moins un coefficient au tableau ; le reste suit.
\end{notepourlaclasse}
